\documentclass[11pt]{article}
\usepackage[utf8]{inputenc}
\usepackage[T1]{fontenc}
\usepackage{lmodern}
\usepackage{enumitem}
\usepackage{amsmath,amssymb}
\usepackage{geometry}
\usepackage{xcolor}
\usepackage{listings}
\usepackage{hyperref}
\geometry{margin=1in}
\setlist[itemize]{leftmargin=*}
\lstset{basicstyle=\small\ttfamily, columns=fullflexible, keepspaces=true,
  frame=single, breaklines=true, backgroundcolor=\color{gray!7}}

\title{Improving the Tactics \texttt{ring} and \texttt{field}\\
 }
\author{Thomas Portet}

\begin{document}
\maketitle

\begin{abstract}
This document describes an extension of Rocq's reflexive algebraic tactics. The
main goal is to make \texttt{ring} and \texttt{field} usable on expressions that
contain functions outside the ring or field fragment. The approach rewrites the
arguments of such functions, computes polynomial greatest common divisors with
an OCaml oracle, and exposes the result to Rocq through ELPI foreign calls.
\end{abstract}

\section{LiberAbaci and the motivation}
LiberAbaci uses Rocq for early mathematics education. Its goals are:
\begin{itemize}[leftmargin=*]
    \item choose a simple ambient theory and number type;
    \item import theorems from other theories, possibly via Trocq;
    \item support declarative proofs with \texttt{assert}, \texttt{enough}, and
      a Lean-inspired \texttt{calc};
    \item favor simple tactics over long rewrite scripts.
\end{itemize}
The supporting tactics are intended to include \texttt{ring}, \texttt{field}, and
\texttt{nra}. Recent work by A. Dehorey also points in the same direction. A
student should, for instance, be able to write an equality such as
\[
 \ln(\pi^2+\pi\times\pi)=\ln(2\pi^2)
\]
and let the tactic justify the algebraic step inside the logarithm, without a
sequence of manual rewrites.

\section{Where reflexive tactics are insufficient}
Consider
\[
  \frac{\pi}{\pi^2+\pi^2}=\frac{1}{2\pi}.
\]
\texttt{field} can prove this if it proves $\pi^2+\pi^2\ne0$ and
$\pi\ne0$.
In practice,\texttt{field} reduces the equality to these non-zero side conditions; they still have to be proven separately, for example using the positivity of
$\pi$.
Now consider the variant
\[
  \log\!\left(\frac{\pi}{\pi^2+\pi^2}\right)
  = \log\!\left(\frac{1}{2\pi}\right).
\]
Here \texttt{field} fails because $x\mapsto\log(x)$ is not a field
expression. The same issue occurs for any outer function that is opaque to the
tactic, such as $\exp$, $\sin$, or $\cos$.
The same issue appears with \texttt{ring}:
\[
  \pi^2+\pi\mathbin{\times}\pi=2\pi^2
\]
is easy for \texttt{ring}, but
\[
  \ln(\pi^2+\pi\mathbin{\times}\pi)=\ln(2\pi^2)
\]
looks like an equality between two different arguments to the tactic.

\section{How \texttt{ring} works}
\texttt{ring} works with integer constants, multiplication, addition,
subtraction, and variables. Other expressions are treated as opaque variables:
\begin{enumerate}
  \item reify the formula by traversing its syntax and replacing each supported
    operation by a constructor in an internal polynomial expression; and replacing each expression starting with an unsupported operation by a fresh variable.
  \item normalize the internal expression by collecting coefficients and
    combining similar monomials;
  \item conclude when the normalized expressions are syntactically equal, since
    the normalization procedure has proved that they denote the same ring term.
\end{enumerate}
For example, replacing $\pi$ by $X$ gives
\[
  \pi(\pi\times1)+\pi\times\pi
  \quad\longmapsto\quad X(X\times1)+X\times X.
\]
This normalizes to $2X^2$ without any knowledge of $\pi$. But
\[
  \ln(\pi(\pi\times1)+\pi\times\pi),\qquad \ln(2\pi^2)
\]
become two distinct variables to the tactic, which cannot see that their
arguments are equal.

\section{Deep ring simplification}
The idea is to simplify function arguments repeatedly. We collect the relevant
subexpressions, apply ring simplification, and repeat until no further progress
is made; then the proof closes by reflexivity.
We only call \texttt{ring\_simplify} on the relevant subexpressions, namely those
that occur as arguments of otherwise opaque function applications.
For example,
\[
 \cos(x+\cos(x+x))
 =\cos(\cos(2x)+(1-(1-x))),
\]
the relevant subexpressions are $x+\cos(x+x)$, $x+x$,
$\cos(2x)+(1-(1-x))$, and $2x$. Simplification turns $x+x$ into $2x$ and
$\cos(2x)+(1-(1-x))$ into $x+\cos(2x)$. The equality becomes
\[
  \cos(x+\cos(2x))=\cos(x+\cos(2x)).
\]
Reification treats a subterm outside the supported ring language as an atom.
Thus, at one level, $x+\cos(x+x)$ is represented as the polynomial
$X_1+X_2$, where $X_1$ represents $x$ and $X_2$ represents
$\cos(x+x)$. Deep simplification first visits the argument $x+x$, normalizes
it to $2x$, and then reifies the outer expression again. The atom for the
function application can consequently be rebuilt as $\cos(2x)$, allowing the
same process to continue at enclosing levels.

\section{The remaining problem with \texttt{field\_simplify}}
The standard \texttt{field} tactic uses fractions as normal forms, but it compares
those fractions by cross multiplication:
\[
  q\ne0\ \land\ s\ne0
  \quad\Longrightarrow\quad
  \left(\frac{p}{q}=\frac{r}{s}
  \quad\Longleftrightarrow\quad
  p\times s=q\times r\right).
\]
This means a fraction does not need to be reduced in order to prove an equality.
For example, $3\pi/6=2\pi/4$ is established by checking the cross-product
identity $3\pi\times4=6\times2\pi$ together with the non-zero conditions on the
denominators.
The difficulty is subtler when the fraction occurs inside a function. The
tactic can prove that

$$
  \frac{3\pi}{6}=\frac{2\pi}{4}
$$

by cross multiplication, but this proof does not require either fraction to
be rewritten to a common representative. In particular, the two sides can
remain syntactically different after field normalization.

This distinction does not matter when the equality of the two fractions is
the final goal. Once the cross-product equality and the required non-zero
conditions have been established, the proof is complete. It matters when the
fraction occurs as an argument of a function that is opaque to the field
normalizer. Consider

$$
  \cos\left(\frac{3\pi}{6}\right)
  =
  \cos\left(\frac{2\pi}{4}\right).
$$

The field normalizer can reason about the two fractions, but it does not
normalize the surrounding occurrences of \texttt{cos}. The two function
applications can therefore contain different terms even when their
arguments have been proved equal.

The required operation is consequently stronger than proving equality of
the two fractions. We need to transform both fractions to a common
representative that can be used in the surrounding expression. For the
example above, a suitable representative is

$$
  \frac{\pi}{2}.
$$

The desired transformation is therefore

$$
  \cos\left(\frac{3\pi}{6}\right)
  \longrightarrow
  \cos\left(\frac{\pi}{2}\right)
  \longleftarrow
  \cos\left(\frac{2\pi}{4}\right).
$$

Cross multiplication alone does not provide this representation. It
establishes an equality between fractions, but it does not require the
numerator and denominator of each fraction to be divided by their common
factor. This is the difference between proving that two expressions are
equal and computing a representation that is useful for rewriting.

A polynomial gcd provides the missing operation. Given a numerator
$N$ and denominator $D$, we compute a common factor $G$ and corresponding
quotients. The numerator and denominator can then be represented by
reduced terms. For example,

$$
  N=X^2-1,\qquad D=X+1
$$

have the common factor $G=X+1$, giving

$$
  N=(X-1)G,\qquad D=1G.
$$

The fraction can therefore be represented by

$$
  \frac{X-1}{1}.
$$

The reduction is useful here because the resulting numerator and
denominator can be used as a common representation when the fraction occurs
below an opaque function. The gcd computation is not needed to establish
the equality of two fractions by cross multiplication. It is needed to
produce a representation that can be propagated through the surrounding
term.

This distinction also explains why the extension is useful for the
reflexive tactics considered here. The normalizer already has enough
information to establish many equalities. The additional computation is
used when the normal form itself must be exposed inside a larger expression.


\section{Common structure for \texttt{field} and \texttt{ring}}
Both tactics normalize terms and then check equality. For \texttt{ring}, the
normal forms are polynomials and equality is syntactic. For \texttt{field}, the
normal forms are fractions. Before cross multiplication can be used, the tactic
records the non-zero denominator conditions; the equality itself is then reduced
to a polynomial identity. For instance, proving
\[
  \frac{\pi}{\pi^2+\pi^2}=\frac{1}{2\pi}
\]
requires the cross-product equation together with proofs that
$\pi^2+\pi^2\ne0$ and $\pi\ne0$.

\section{The gcd oracle}
The field-simplification theorem has the following shape:
\[
 \operatorname{Fnorm}(l,e)=r \Longrightarrow
 \operatorname{FEeval}(l,e)=
 \frac{\operatorname{Pphi\_pow}(l,\operatorname{norm}(\operatorname{num}(r)))}
  {\operatorname{Pphi\_pow}(l,\operatorname{norm}(\operatorname{den}(r)))}.
\]
The purpose of the oracle is to compute a common factor together with the
quotients required by the field-simplification theorem. For a simple
non-trivial example, consider

$$
  N=X^2-1=(X-1)(X+1),
  \qquad
  D=X+1.
$$

The common factor is

$$
  \gamma=X+1.
$$

The corresponding quotients are

$$
  \nu=X-1,
  \qquad
  \delta=1.
$$

The reduced fraction is therefore represented by

$$
  \frac{\nu}{\delta}
  =
  \frac{X-1}{1}.
$$

The oracle does not directly return an equality proof. Instead, it returns
the data from which Rocq can construct such a proof. More precisely, the
extension introduces a non-zero integer $m$ and requires

$$
  mN=\nu\gamma,
  \qquad
  mD=\delta\gamma.
$$

The integer $m$ accounts for the fact that the polynomial representation
used by the oracle may differ from the normalized representation used by
Rocq by a non-zero scalar.
The scale factor is therefore part of the interface between the two
polynomial representations. The oracle is free to return a polynomial
representation that is equivalent up to a non-zero integer factor, while the
Rocq side can account for this factor explicitly in the factorization
identities.

The two identities
\[
  mN=\nu\gamma,
  \qquad
  mD=\delta\gamma
\]
show that the same factor $\gamma$ has been removed from the numerator and
denominator. Since $m\ne0$, this scaling does not change the represented
fraction. The field theorem can consequently replace the original numerator
and denominator by $\nu$ and $\delta$.
The Rocq side subsequently proves these two factorization identities and
the condition

$$
  m\ne0.
$$

The output of the oracle is therefore treated as computational data, while
the identities that justify its use are established inside Rocq.

The resulting theorem uses the reduced numerator and denominator:

$$
  \operatorname{FEeval}(l,e)=
  \frac{\operatorname{Pphi\_pow}(l,\operatorname{norm}(\nu))}
       {\operatorname{Pphi\_pow}(l,\operatorname{norm}(\delta))}.
$$

Thus the oracle supplies a candidate reduced representation, and the proof
script verifies the algebraic relations required to use it.

This reduction should be distinguished from the usual proof of equality of
fractions. Cross multiplication is sufficient for such a proof. The gcd
computation instead provides numerator and denominator quotients that can be
used to construct a common representation. This is the representation
needed when the fraction occurs below an otherwise opaque function.
\section{A complete example}

Consider a fraction occurring inside an opaque function,

$$
  f\left(\frac{X^2-1}{X+1}\right),
$$

where the function $f$ is not part of the polynomial language handled by
the field normalizer. The field normalization procedure can reason about
the fraction, but it does not by itself require the fraction to be reduced
to a common representative.

The numerator and denominator are

$$
  N=X^2-1,\qquad D=X+1.
$$

The polynomial gcd computation returns

$$
  \gamma=X+1,
  \qquad
  \nu=X-1,
  \qquad
  \delta=1.
$$

The returned data satisfy

$$
  N=\nu\gamma,
  \qquad
  D=\delta\gamma.
$$

In this example the scale factor is $m=1$.

The fraction can consequently be represented as

$$
  \frac{N}{D}
  =
  \frac{\nu}{\delta}
  =
  \frac{X-1}{1}.
$$

The useful result is not only the proof that the two fractions are equal.
The reduced numerator and denominator provide a new representation that can
be inserted into the surrounding expression:

$$
  f\left(\frac{X^2-1}{X+1}\right)
  \longrightarrow
  f\left(\frac{X-1}{1}\right).
$$

The example illustrates the role of the gcd computation in the extension.
The polynomial gcd is used to obtain the common factor that is removed from
the numerator and denominator. Rocq then verifies the factorization
identities returned by the computation. The external computation therefore
determines the candidate representation, while the proof script establishes
the equalities needed to justify the rewrite.

\section{The polynomial oracle in more detail}
The OCaml implementation does not manipulate the input expressions directly
when computing the gcd. It first translates them into a sparse multivariate
representation. A monomial consists of a rational coefficient and a sorted list
of pairs $(i,e)$, where $i$ is a variable index and $e$ is its exponent. A
polynomial is a list of such monomials. Addition collects monomials with the
same exponent vector, while multiplication merges exponent vectors and
multiplies coefficients.

Rational coefficients are important because normalization can introduce
non-integral leading coefficients. The implementation stores them explicitly as
pairs of integers and reduces them after arithmetic operations. The surface
syntax remains deliberately small:

\begin{lstlisting}[language=ML]
type rat = { num : int; den : int }

type poly =
  | Const of rat
  | Var of int
  | Add of poly * poly
  | Mul of poly * poly
\end{lstlisting}
For example, the monomial
\[
  3X_1^2X_4
\]
is represented by the coefficient $3$ together with the exponent list
\[
  [(1,2),(4,1)].
\]
A polynomial is then a collection of such monomials. Terms with the same
exponent list are combined during addition, while multiplication adds the
corresponding exponents and multiplies the coefficients.

This representation is sparse: variables that do not occur in a monomial
do not appear in its exponent list. This avoids representing a polynomial
as a dense array whose size would depend on the largest variable index or
degree.
After conversion to the sparse representation, monomials are ordered
lexicographically, with lower variable indices taking precedence. This order is
used by leading-monomial selection, exact monomial division, polynomial
reduction, and Buchberger's algorithm. The latter constructs a Groebner basis
from the S-polynomials of the current basis elements and reduces each new
remainder until all pairs have been processed.
The ordering is also used to select a leading monomial during polynomial
reduction. Given a polynomial and a basis element, the implementation first
checks whether the leading monomial of the basis element divides the leading
monomial of the polynomial being reduced. If so, the corresponding multiple
of the basis element is subtracted. Otherwise the leading monomial is moved
to the remainder. Repeating this process gives the normal remainder with
respect to the current basis.

Buchberger's algorithm applies this reduction to the $S$-polynomials of pairs
of basis elements. The resulting basis is then used by the gcd computation
described above.
The gcd computation uses an auxiliary variable $t$. For input polynomials $p$
and $q$, it considers the two generators
\[
  t p \qquad\text{and}\qquad (1-t)q.
\]
The $t$-free part of the ideal generated by these polynomials contains a
generator for the least common multiple of $p$ and $q$. The implementation
therefore shifts the original variables by one position, reserves variable $0$
for $t$, computes a Groebner basis, and extracts the basis elements that do not
contain $t$. After interreduction, the resulting LCM is used to obtain
\[
  \gcd(p,q)=\frac{p q}{\operatorname{lcm}(p,q)}.
\]
\section{Two polynomial representations}
Rocq uses \texttt{PExpr} and \texttt{Pol}. The function \texttt{norm} converts
\texttt{PExpr} into the sparse Horner form \texttt{Pol}. ELPI uses
\texttt{polyT}. The bridge converts between these representations.


For example, the polynomial

$$
  X_1+X_1^2
$$

is represented in \texttt{PExpr} as

$$
  \texttt{PEadd (PEX 1) (PEmul (PEX 1) (PEX 1))}.
$$

Applying \texttt{norm} produces its canonical \texttt{Pol} representation.
The corresponding ELPI expression,

$$
  \texttt{add (var 1) (mul (var 1) (var 1))},
$$

is decoded into this \texttt{PExpr} representation before normalization.
The reverse conversion is used to return the result of the polynomial
computation to ELPI.

This gives the pipeline three distinct representations: ELPI terms are
convenient for pattern matching and foreign calls, \texttt{PExpr} is the
Rocq syntax accepted by the normalizer, and \texttt{Pol} is the canonical
sparse representation used by the field lemmas. Keeping these
representations separate lets the OCaml code remain independent of Rocq's
internal polynomial constructors.

\texttt{Pol} representation.
% something done right ring normalization mahboubi TODO CITATION 
For example, an ELPI expression such as
\texttt{add (var 1) (mul (var 1) (var 1))} is first decoded to a Rocq
\texttt{PExpr}, corresponding to $X_1+X_1^2$. Rocq then applies
\texttt{norm} to obtain its sparse Horner representation in \texttt{Pol}. The
reverse direction is used for the result returned by the oracle. This gives the
pipeline three distinct representations: ELPI terms are convenient for pattern
matching and foreign calls, \texttt{PExpr} is the Rocq syntax accepted by the
normalizer, and \texttt{Pol} is the canonical sparse representation used by
the field lemmas. Keeping these representations separate lets the OCaml code
remain independent of Rocq's internal polynomial constructors.
\section{ELPI as the bridge}
ELPI is a Prolog-like language for term manipulation. Its foreign-function
interface is implemented in OCaml and connects ELPI to OCaml code loaded by the
Rocq plugin. The plugin exposes the polynomial representation to ELPI through
constructors for constants, variables, addition, and multiplication.

The plugin exposes \texttt{polyT} with constructors \texttt{pcst},
\texttt{var}, \texttt{add}, and \texttt{mul}:
\begin{lstlisting}[language=ML]
K("add", "binary addition", S(S(N)),
  B (fun x y -> Add (x, y)),
  M (fun ~ok ~ko t ->
    match t with Add (x, y) -> ok x y | _ -> ko ()))
\end{lstlisting}
Here \texttt{S} describes the recursive value and \texttt{ko} handles a
mismatch.

\section{Encoding and foreign calls}
The decoder builds a \texttt{PExpr} term and asks Rocq's VM for its
\texttt{Pol} form:
\begin{lstlisting}
func pol_decode polyT -> term.
pol_decode PolyT_a Result :-
  pe_decode_naive PolyT_a Pe_a,
  coq.reduction.vm.norm {{norm lp:Pe_a}} {{Pol Z}} Result.
\end{lstlisting}
OCaml provides
\[
  \texttt{Gcd.poly\_gcd : poly -> poly -> (int * poly * poly * poly)}.
\]
This becomes an ELPI predicate with two inputs and four outputs:
\begin{lstlisting}
kind polyT type.
external symbol pcst : int -> polyT.
external symbol var  : int -> polyT.
external symbol add  : polyT -> polyT -> polyT.
external symbol mul  : polyT -> polyT -> polyT.
external func gcd_poly polyT, polyT -> int, polyT, polyT, polyT.
\end{lstlisting}
The wrapper encodes the inputs, calls the function, and decodes the results:
\begin{lstlisting}
func gcd_and_factors term, term -> term, term, term, term.
gcd_and_factors N D Ne De Gcd MZ :-
  pol_encode N Nx, pol_encode D Dx,
  gcd_poly Nx Dx M G N' D',
  pol_decode N' Ne, pol_decode D' De,
  pol_decode G Gcd, z_decode M MZ.
\end{lstlisting}

\section{Ltac calling ELPI}
The proof script still lives in Ltac, but the actual computation is delegated to
ELPI. The bridge is explicit: an Ltac function builds a Rocq term, calls the
ELPI computation, and receives back the reduced numerator, denominator, and gcd:
The value named \texttt{res} is a nested Rocq pair returned by the ELPI tactic.
The first projection is the integer scale factor, while the remaining
projections contain the reduced numerator, reduced denominator, and gcd. The
following calls to \texttt{cbv} only unpack this transport structure; they do
not perform the polynomial computation again. The later assertions connect the
oracle output back to Rocq by proving the two factorization identities and the
non-zero condition on the scale factor.
\begin{lstlisting}
let res :=
  constr:(ltac:(elpi factorize_by_gcd ltac_term:(N1) ltac_term:(D1))) in
let F := eval cbv [fst] in (fst res) in
let N2 := eval cbv [fst snd] in (fst (snd res)) in
let D2 := eval cbv [fst snd] in (fst (snd (snd res))) in
let Gcd := eval cbv [snd fst] in (snd (snd (snd res))) in
let fact_n0 := fresh "factor_not_0" in
let num_eq := fresh "equality_for_numerator" in
let den_eq := fresh "equality_for_denominator" in
assert (fact_n0 : IZR F <> 0) by
  (apply eq_IZR_contrapositive; easy);
enough
  ((Pmul (Pc F) N1 ?== Pmul N2 Gcd) = true /\
   (Pmul (Pc F) D1 ?== Pmul D2 Gcd) = true)
  as [num_eq den_eq];
  exact (hyp F N2 D2 Gcd fact_n0 num_eq den_eq).
\end{lstlisting}
This is an Ltac function rather than a standalone tactic: it takes terms and
proof data as arguments and performs the final proof steps. The cross-product
equalities are proved only after the oracle has returned its factors, and the
resulting hypothesis is then used to rewrite the original term. Ltac chooses the
terms to send, ELPI encodes them and runs the polynomial routine, and the tuple
coming back is consumed by the proof script.
\section{Architecture of the extension}

The extension separates the proof-producing part of the tactic from the
polynomial computation. Ltac remains responsible for interacting with the
existing field-simplification infrastructure and for constructing the final
proof. ELPI is used as the interface between this Ltac code and the
polynomial computation implemented in OCaml.

The data passes through several representations. The original numerator and
denominator are Rocq terms. They are passed to ELPI, where they are encoded
as values of type \texttt{polyT}. These values provide a representation that
can be inspected and passed to foreign functions. The decoder then converts
them to Rocq's \texttt{PExpr} representation.

The \texttt{PExpr} representation is useful because it is accepted by
Rocq's polynomial normalizer. Applying \texttt{norm} converts the expression
to the corresponding \texttt{Pol} representation. The polynomial data can
then be passed to the OCaml gcd implementation.

The reverse path is used for the result. The polynomial factors returned by
the OCaml computation are converted back to \texttt{polyT}, decoded to
Rocq terms, and normalized to \texttt{Pol} when necessary. ELPI returns the
result to the Ltac code as a tuple containing the scale factor, reduced
numerator, reduced denominator, and gcd.

The separation between these representations has a practical purpose. ELPI
terms provide a convenient interface for pattern matching and foreign
calls, while \texttt{PExpr} and \texttt{Pol} remain the representations used
by Rocq's existing polynomial machinery. The OCaml implementation therefore
does not need to construct or manipulate Rocq's internal polynomial
representations directly.
The complete computation can be summarized as follows:
\[
\begin{array}{c}
\text{Rocq term}
\\
\downarrow
\\
\texttt{polyT}
\\
\downarrow
\\
\texttt{PExpr}
\\
\ \ \    \ \ \   \mathord{\downarrow}\texttt{norm}
\\
\texttt{Pol}
\\
\downarrow
\\
\text{OCaml gcd computation}
\\
\downarrow
\\
\text{resulting polynomial data}
\end{array}
\]
The result then follows the reverse path before being consumed by the Ltac
proof script.

This architecture also separates computation from proof construction. The
OCaml code computes polynomial data, but the Ltac code subsequently uses
this data to construct the equalities required by the field theorem. The
external computation therefore supplies information used by the proof
rather than replacing the proof of the algebraic identities.

\section{Making the computation available}
The extension provides an ELPI entry point used by the Ltac function. It computes
the factors, normalizes them in Rocq, and returns the least common multiple,
reduced numerator, reduced denominator, and gcd:
\begin{lstlisting}
Elpi Tactic factorize_by_gcd.
Elpi Accumulate Plugin "ext.elpi".
Elpi Accumulate File encode.
Elpi Accumulate lp:{{
  solve (goal _ _ _ _ [trm N, trm D] as G) GL :-
    gcd_and_factors N D N' D' Gcd LCM,
    coq.reduction.vm.norm {{norm lp:N'}} {{Pol Z}} N'',
    coq.reduction.vm.norm {{norm lp:D'}} {{Pol Z}} D'',
    refine {{pair lp:LCM (pair lp:N'' (pair lp:D'' lp:Gcd'))}} G GL.
}}.

\end{lstlisting}
%add stuff there
\section{Integrating with \texttt{field}}
The workaround is a parameterizable Ltac function with three parameters: a
normalization function, a justification theorem, and a finishing tactic. The
existing field code computes unknowns, applies a lemma, and rewrites. The
extension collects the numerator and denominator, invokes the gcd computation,
and rewrites using the resulting data. The modified pieces include versions of
\texttt{Field simplify gcd}, \texttt{Field norm gen},
\texttt{ReflexiveRewriteTactic}, and \texttt{ApplyLemmaThenAndCont}.

In a typical call, the function receives the term being rewritten, the
hypothesis produced by field normalization, and the numerator and denominator:
\begin{lstlisting}[language=ML]
fraction_finisher Term hyp D N.
\end{lstlisting}
It then calls \texttt{factorize\_by\_gcd} to obtain the reduced data, proves the side conditions and applies the resulting equality. The ELPI entry point is
therefore an auxiliary computation used by Ltac rather than the user-facing
tactic itself.

\section{Lessons learned}
This experiment brings together four languages: Rocq,
Ltac, ELPI, and OCaml. Ltac remains the central layer. A direct translation to
OCaml is possible, but it is much heavier; ELPI is shorter and has access to
both type checking and conversion. The Ltac infrastructure is powerful but not
always easy to read or debug. 
The current extension addresses the representation problem for a field
expression selected by the existing field-simplification infrastructure. It
does not yet provide the same recursive traversal of arbitrary nested field
expressions as the deep ring procedure described above. A more general deep
field simplification strategy is therefore left for future work.
\section{Future work}
Future work includes deep field simplification, more careful staging, fewer
generated conditions, support for exponents in deep ring and field reasoning,
commutative-group isomorphisms, and a stronger polynomial gcd computation with
unbound integers.

\section{Conclusion}
Ring normalization can simplify expressions even when they appear under function
symbols. An OCaml gcd oracle can reduce field fractions to a common shape, and
ELPI provides the bridge between the proof assistant, the tactics, and the
external computation.

\end{document}
